\clearpage 
\section{Coursework 1 (Part 1): \\ Simple Robots}

The first coursework for CS118 consists of two parts. You will need to 
complete both parts if you are aiming for a top grade. \\

Your answers to this coursework must be completed by \deadlineone. 
You will be asked to come in on that day and demonstrate 
that your code works and that you have understood the material for each of
the exercises. If we are not able to mark your work on that day, for whatever 
reason, then you will receive no marks.\\

Your work will of course be marked for functionality, that is the program does 
what it is supposed to do. It is also useful to remember that the work will 
be marked for programming style, clarity, re-usability and use of 
techniques taught throughout the course. This means that even if your code 
works wonderfully, you may not get fantastic marks if it looks like a 
dog's dinner. Likewise, if you do not finish all the exercises, but your
solutions look like a masterpiece, then you are likely to do well. \\

{\bf Cooperation, Collaboration and Cheating}\\

If a submitted program is not entirely your own work, you will be required
to state this when the work is marked. Any and all 
collaboration between students
must be acknowledged, and may result in stricter marking of the work.
Consultation of textbooks is encouraged, but programs described elsewhere
should not be submitted as your own, even if alterations are made.
It will be useful to quote here the University's regulations on the subject:

\begin{quote}
{\em \ldots `cheating' means an attempt to benefit oneself, or another, by
deceit or fraud. This shall include deliberately reproducing the work of
another person or persons without acknowledgement. A significant amount of
unacknowledged copying shall be deemed to constitute prima facie evidence of
deliberation, and in such cases the burden of establishing otherwise shall
rest with the candidate against whom the allegation is made.}
\end{quote}

\noindent
Therefore, it is as serious for a student to permit work to be copied as it
is to copy work. Any assistance you receive must be acknowledged.
If in doubt, ask. It should also go without saying that you are prohibited 
from uploading your solutions to the internet.

\subsection{Getting started}

To begin you need to copy the robot-maze environment and 
the controller software to your home directory. I suggest that you first 
create a {\tt cs118} directory. Invoke a terminal window and type in 
this window the UNIX command \\

{\tt mkdir cs118} \\

\noindent
and then change to this directory using the command {\tt cd cs118}. \\

\noindent
You should now use a web browser to go to the CS118 course web page 
(which should be in your bookmarks if you have followed all the instructions
in Chapter 2). \\

\noindent 
Under the {\bf Coursework} section of this web page you will see that there 
are four links, one of which says {\bf Maze software} and one of which says 
{\bf Dumbo controller}. Click on these with your right mouse button and 
select the {\em Save Link Target As...} option from the menu. This will allow you 
to save the file. When you save these files, make sure you double-click on the 
{\tt cs118} directory so that the file is saved to the appropriate place.
The {\bf Download Manager} will confirm that the file has been saved and 
once you have done this for both files, you can check that the files have 
been saved by typing {\tt ls -l} in your terminal window (you may need to navigate
to the correct folder using some of the commands you encountered earlier).\\

\noindent
You should find two new files in this directory. 
One of these files is a {\tt .jar} file;  
this postfix means that the file is a Java archive, an 
efficient ZIP-like file format which allows all the component parts of the 
robot-maze environment (stored in this file) to take up as little space as 
possible in your home directory. The other file is a {\tt .java} file 
like those you created in your first seminar. This {\tt .java} file is the 
robot controller which interfaces with the robot-maze environment software. \\

\noindent
The result of the {\tt ls -al} command should look something like this:

\begin{verbatim}
-rw-------   1 saw   dcsstaff     697 Aug 11 10:32 DumboController.java
-rw-------   1 saw   dcsstaff   99539 Aug 11 10:32 maze-environment.jar
\end{verbatim}

\noindent
If you find that the file size of either of these files is zero
(this is the number in the fifth column), then something has gone wrong during the 
downloading. In this case you should try downloading them once again. \\

\noindent
You need to compile the {\tt .java} file if you are to run it and in order for
the controller software to run along side the robot-maze environment
the two programs need to be compiled together. This requires a small 
addition to the {\tt javac} command which you used in compiling your 
first Java programs. Type into you terminal window the command \\

{\tt javac -classpath maze-environment.jar DumboController.java} \\

\noindent
This compiles the controller program {\tt DumboController.java}
into a corresponding class file ({\tt DumboController.class}). The 
compilation is performed in the context of the 
robot-maze environment (through the {\tt -classpath maze-environment.jar} 
extension) which is just what we want.\\

\noindent
You can now run the robot-maze environment by typing \\

{\tt java -jar maze-environment.jar \&} \\

\noindent
The {\tt \&} symbol runs the java command in the background,
such that it frees the window so that you can still use it for 
other business. Admire the baroque elegance of the highly sophisticated
computer graphics in the robot-maze environment program. To change
the maze, click on the {\tt Generators} button and then on 
the {\bf PrimGenerator} in the window above. You will see that this fills
the {\em Current Generator} information panel. If you now click on the 
{\bf New Maze} button at the bottom right you will 
get a new maze (generated through an application of Prim's algorithm). \\

\noindent
Now that you have a maze you need a robot. Click on the {\bf Controllers} 
button
and select the {\bf RandomRobotController}. You will see that this configures
the {\em Current Controller} information panel. \\

\noindent
The {\bf RandomRobotController} is a pre-installed piece of controller 
software which drives the direction-choosing capabilities of a basic robot. 
Before clicking on the {\bf Start} button to test the robot, set the 
{\bf Speed} gauge to the far right of the screen. \\

\noindent
When the robot is running you will see that it exhibits some unusual behaviour.
First you will see the direction change, as indicated by the blue arrow. You 
will also notice that it leaves a trail (of been-before squares) as it moves 
through the maze. Occaisionally the robot crashes into a wall 
(indicated in red), this is because the 
controller which is being used to drive the robot makes no allowance for 
where the walls are in the maze. \\

\noindent
You should familiarise yourself with this environment. See what happens 
when you increase the speed, try generating new mazes, try editing mazes 
with your mouse, try moving the location of the target, 
try changing the dimensions of the maze.

\subsection{Loading controllers}

One of the nice features of the robot-maze environment is 
the ability to experiment with different controller software. The 
{\tt DumboController} which you compiled earlier can be loaded into the 
environment by clicking {\bf Controllers} followed by {\bf Add}.\\

\noindent
This will provide you with a directory menu from which you should double 
click on your {\tt cs118} directory. Here you will find your 
{\tt DumboController.class} file which you can then highlight and {\bf Open}.\\

\noindent
You will see that this adds the {\tt DumboController} to your list of robot
controllers. You will use this same process to load all the robot controllers
which you write during this first piece of coursework.\\

\noindent
If you click on the {\bf DumboController} in the robot controllers menu, the 
environment will switch between controllers. Run the new robot controller
to see if you can work out where {\bf RandomRobotController} and 
{\bf DumboController} differ.  Try to 
characterise the strategies which they appear to follow.  You
may find find it helpful to test the robots on different mazes. \\

\noindent
It may become obvious that it is not always easy to see exactly what
strategies these robot controllers appear to follow. It is often quite 
difficult to reverse engineer from the {\em behaviour} to the 
{\em specification}. Similarly, it is not good practice to have 
a specification-less program, as if the user is in any
doubt as to the behaviour of part of the program, this
problem can be easily resolved by consulting the specification. 

\subsection{Analysing code}

Use a text editor to look at the source code in the file
{\tt DumboController.java}.  The method {\tt controlRobot} is the 
controller part of the code. If you have not already worked it out, study the code and see if 
you can detect what strategy this control program implements. \\

\noindent
Note the use of the {\tt import} statement at the 
top of the program. This is the statement which connects the robot-maze 
environment with the controller code. The behaviour of this interface was
described in Chapter 5 and you might like to go back to this chapter and 
just check what each of the robot methods and constants do. \\

\noindent
Talk with your friends about this code. Is it clear to you which parts of the 
code are `pure' Java and which parts come from the interface to the 
robot-maze environment?

\subsection{Exercise 1}

$\diamondsuit$ An order is received from an existing customer for a modified 
dumbo robot:

\begin{quote}
`Could we have a modified robot controller that still chooses directions
randomly, but avoids crashing into walls.'
\end{quote}

\noindent
This description can be identified as the {\em specification of
requirements} for the new robot which the customer requires. \\

\noindent
A good software developer will set about solving this problem in a 
systematic and logical fashion; for example, using the processes of 
{\it Design}, {\it Build} and {\it Test} which were described in 
Chapter 3.\\

\noindent
Designing a new piece of software requires a complete understanding of 
the problem to hand; you cannot write a program for a problem which you 
do not understand. To help work out what the specification states, we will 
break the description up into its constituent parts. 

\begin{itemize}

\item
The text describes the modified robot as `still choosing directions
randomly'. This would suggest that the part of the robot control program which
chooses a random number and then converts this to a direction should stay as 
it is. 

\item
The text also states that the robot should `avoid crashing into
walls'. Sometimes the existing robot controller chooses a random direction
which will point the
robot towards a wall. What you need to do is filter out these
occurrences. This will involve looking to see if the direction chosen
does point the robot towards a wall, and if so, choosing another
direction. 

\end{itemize}

\noindent
A software developer would be right in thinking that the existing 
{\tt controlRobot} method in the {\tt DumboController.java} file 
can be reused; 
there are many similarities between the existing robot 
controller which you studied in previously and the new robot controller. 
Your answer to this exercise should therefore be based on the code 
found in {\tt DumboController.java}. \\

\noindent
The main difference between the old controller and the new one is that the 
new controller will prevent the robot from crashing into walls. \\

\noindent
{\bf Design question 1:}
How do you think you can detect if the robot is about to crash into a wall? 
Hint: look at Section \ref{lookmethod} of {\it The Guide}. \\ 

\noindent
Once you have discovered how to detect for collisions, you will need to
ensure that the robot controller keeps choosing directions for the 
robot until a non-wall direction is found. \\

\noindent
{\bf Design question 2:}
How do you plan to do this? Hint:
look at the lecture notes, this guide and any Java books you have to find out
more about loops. \\

\noindent
Once you have designed the program you are ready to {\em build} the
new robot controller. 
The new robot controller can be built by making modifications 
to the existing {\tt controlRobot}
method in the {\tt DumboController.java} file. If you've carefully planned
how your new robot should work, it should only require about two
lines of code to be changed or added, so there is no need to get
carried away, or indeed too daunted by the programming task ahead!\\

\noindent
After saving the {\tt DumboController.java} file, you should compile the 
code using the command\\

{\tt javac -classpath maze-environment.jar DumboController.java}\\

\noindent
to generate a new {\tt DumboController.class} file. 
Once you have eliminated any fatal,
compiler-detectable errors, a new class file will be created. 
The new class file will be detected by the robot-maze software and it will 
ask you whether you would like to reload this new class; you should 
press {\bf Yes}; you can now test your new solution. \\

\noindent
Before you finish this exercise, consider how you would convince the
customer that you have tested the program and that it fully meets the
customer's requirements. 

\subsubsection{Performance monitoring}

Now that the robot no longer crashes into walls, it is easier for your
customer to monitor the robot's behaviour. As a result, you receive
an email stating that they have noticed some rather unexpected behaviour. \\

\noindent
While testing the current robot your customer noticed that although
it seems to choose directions randomly, some directions appear to be selected 
more often than others. \\

\noindent
This is a difficult trait to investigate by hand. You could try running your robot 
slowly and making a note of the directions it chooses, but this is rather 
painful (and life is too short for such tedium). An alternative 
approach is to add a logging mechanism, so that the robot keeps a record of which
direction is selected each time it moves. The idea is that from this 
log of movements you will be able to analyse the behaviour of the
robot for a particular maze. \\

\noindent
The following output is taken from a working solution to this exercise:

\begin{verbatim}
I'm going forward at a deadend
I'm going forward down a corridor
I'm going forward down a corridor
I'm going forward at a junction
I'm going backwards at a deadend
I'm going backwards at a junction
I'm going backwards at a deadend
I'm going forward at a junction
I'm going backwards down a corridor
I'm going forward at a junction
I'm going backwards at a deadend
I'm going forward at a junction
I'm going forward down a corridor
I'm going forward down a corridor
I'm going backwards at a deadend
I'm going forward down a corridor
I'm going forward down a corridor
I'm going forward at a junction
I'm going backwards at a deadend
I'm going right at a junction
I'm going forward down a corridor
I'm going right at a junction                                                  
\end{verbatim}

\noindent
You will see that the first thing that the robot detects (as it begins the 
exploration of the maze) is that it is at a dead-end and therefore the only 
thing that it can do in this case is move forward (see Figure~\ref{ex7fig}).\\

\noindent
Once it has done this, it then detects that it is in a corridor and therefore
it decides to move forward. The same thing occurs for the third move.\\

\noindent
After three moves the robot finds itself at a junction, here it decides to go
forward once more, at which point it reaches a dead-end and can go nowhere but
backwards. \\

\begin{figure}
\centering
\includegraphics[width=4cm]{ex7fig}
\caption{Coverage through an example maze.\label{ex7fig}}
\end{figure}

\noindent
You can follow the route of the robot until the reset button was pressed. You 
will see that the robot is not particularly efficient (as it is operating in 
just the same way as the robot in the previous task) but it does
print out an accurate log which itself might be useful at a later date.\\

\noindent
It is always simpler to write more complicated software such as this as a 
series of smaller tasks. We will be building on the results of our previous
exercises, so make sure that any programming which you do is done in the file 
{\tt DumboController.java}. \\ 

\noindent
Consider the first problem of outputting the text which identifies the 
direction in which the robot is heading, the \\

{\tt I'm going forward} \\

\noindent
part of the output (again, refer back to Section \ref{lookmethod} of {\it The Guide}!). \\

\noindent
For the robot controller to print a log of this 
direction chosen, you need to include an instruction to output
text. You have already seen such an instruction in the exercises in 
your first problem sheet. 
You might want to remind yourself how the
{\tt System.out.println()} instruction works. Try adding a simple 
{\tt System.out.println()}
instruction to your robot controller, one that says ``{\tt I'm going }'',
or something similar. Compile and run the new robot controller to
observe its effect. \\ 

\noindent
Now that your controller outputs some text, you are ready to modify the program
so that it outputs the {\tt forward}, {\tt backwards}, {\tt left} or 
{\tt right} as
required. Producing this output is quite simple using the 
{\tt System.out.println} instruction; the trick is deciding which of 
{\tt System.out.println("forward");} or 
{\tt System.out.println("backwards");} etc. is required. \\

\noindent
Let's consider the sub-problem of recording the direction chosen.
We could use the {\tt System.out.print()} to build our output string as we go,
or we could store {\tt "forward"}, {\tt "backwards"}, etc. in a string and combine
all variables later. We could even query the direction variable to find which direction
we have chosen to face and use a conditional statement (hint: such as a {\tt switch})
to write out the correct part of the sentence. The main point here, is that there are
many ways to achieve the same goal, try a few and pick your favourite.\\

\noindent
$\diamondsuit$ Design a program which prints the direction the robot is going in.\\

\noindent
Once you have designed this part of the program, it may be worth 
building it and running some tests to make sure that things are working
correctly. \\

\noindent
Now you are ready to work on the final part of the software; detecting 
whether the robot is at a dead-end, in a corridor, at a junction or at a 
crossroads. Answer the 
following questions as part of your design: \\


\noindent
{\bf Design question 1:}
How is it possible for the robot to detect whether it is at a 
dead-end, in a corridor, at a junction or at a 
crossroads? \\

\noindent
{\bf Design question 2:}
Is there a corresponding method in the maze-environment package which allows 
this detection to take place?
Hint: see Section 5.2.3. \\

\noindent
{\bf Design hint 1:}
You might like to add an additional variable to you code, {\tt walls} say, in
which you store the number of walls surrounding the robot. Once you have 
done this you will want to answer the following design questions: \\

\noindent
{\bf Design question 3:}
Where in the control program do you want to declare this variable? Think carefully
about the required {\bf scope} of the variable.\\

\noindent
{\bf Design question 4:}
Where in the control program do you initialise this variable and to what 
value is it initially set?\\

\noindent
{\bf Design question 5:}
What criteria must be fulfilled if values are to be assigned to this 
variable?\\

\noindent
{\bf Design question 6:}
Where in the code is this value inspected and what is the result of this?\\

\noindent
This should allow you to establish a design for the second part of this 
program. \\

\noindent
$\diamondsuit$ Build the
code and design some test criteria. You might want to look at how you 
would test the logical paths of your code -- this was previously 
discussed in Section 3.3 of {\it The Guide}. \\

\noindent
It is worth noting that the additional code should be about twenty lines 
long. \\

\noindent
Once you have thoroughly tested your solution so that you are happy that it meets the 
required behaviour\footnote{paying careful attention to the layout of the logging output},
copy your solution to {\tt Ex1.java} (Make sure you change the class name and check your code
compiles {\bf after} you've done this).

\subsubsection{Assessing performance}

$\diamondsuit$
Using the logging from the previous exercise, determine whether the customer was 
correct in stating that the robot chooses some directions more often than others.\\

\noindent
{\bf Hint:}
Notice how often the robot goes {\tt LEFT} compared to how often it goes {\tt RIGHT}, and 
how often the robot chooses to go {\tt AHEAD} compared to how often it goes {\tt BEHIND}. 
You may find it helpful to run the controller on different mazes. Perhaps even a blank maze?\\

\noindent
Investigating the bias in choice of directions is difficult by hand. However, 
what we can do is use some additional code to analyse the log which the robot
prints out.\\
  
\noindent
We have been supplied with some analysis software (by our sceptical customer).
Provided that your logging output from the previous exercise is correct, this
software will count the number of times the robot heads in each of the four 
directions. \\

\noindent
You should download this analysis software from the course web page. Here 
you will find a link to the file called {\tt count.pl}. Right click on this 
file and {\em Save Link Target As} to save the file in your {\tt cs118} 
directory, where you have stored your answers to the exercises and the 
maze environment. Check that the file has been downloaded properly by typing
the command {\tt ls -al} in a terminal window and inspecting the size 
of the file. \\

\noindent
The {\tt count} program will examine your output from Exercise 5 and produce
a summary of the moves of the robot. For this reason, your output from 
Exercise 5 must be precisely formatted, otherwise the analysis will not 
work properly. To run the maze environment alongside the {\tt count} program 
you should type \\

{\tt java -jar maze-environment.jar | ./count.pl \&}\\

\noindent
If you are using your own computer for this assignment, with the Windows operating
system, you will need to download a different count application, and it is used in a slightly 
different way. Instruction for this are posted on the course web-page.\\

\noindent
The {\tt | ./count.pl} part of the command tells the computer
to send your output to the {\tt count} program to be analysed.\\

\noindent
If you have done everything correctly\footnote{If, for some reason, this 
does not work, do two things: 1. check that the file you have saved is called
{\tt count.pl} and if it is not do a {\tt mv x count.pl}, where {\tt x} is the 
name of the file you saved; 2. make sure that it executes correctly, you 
can ensure this by typing {\tt chmod u+x count.pl} in the command window.}
then the robot
should exhibit the same behaviour as before, but will output some different
information on the console. This will be the analyser's output, and
will look something like:

\begin{verbatim}
Summary of moves: Forward=1 Left=0 Right=0 Backwards=0
Summary of moves: Forward=2 Left=0 Right=0 Backwards=0
Summary of moves: Forward=3 Left=0 Right=0 Backwards=0
Summary of moves: Forward=4 Left=0 Right=0 Backwards=0
Summary of moves: Forward=4 Left=0 Right=0 Backwards=1
Summary of moves: Forward=4 Left=0 Right=0 Backwards=2
Summary of moves: Forward=4 Left=0 Right=0 Backwards=3
Summary of moves: Forward=5 Left=0 Right=0 Backwards=3
Summary of moves: Forward=5 Left=0 Right=0 Backwards=4
Summary of moves: Forward=6 Left=0 Right=0 Backwards=4
Summary of moves: Forward=7 Left=0 Right=0 Backwards=4
Summary of moves: Forward=8 Left=0 Right=0 Backwards=4
Summary of moves: Forward=9 Left=0 Right=0 Backwards=4
\end{verbatim}

\noindent
Examine whether the summary is actually consistent with the robot's moves. 
If you find that the output seems wrong, it is most likely that the log of
movements you printed for the previous exercise 
has not been formatted correctly. Perhaps you have
forgotten to add a space character, or your output has a typing error...\\

\noindent
There is a lesson here in always paying careful attention to the {\em 
program specification}. It might appear a minor problem to forget a space 
in a log of robot movements, or you might decide that the output looks nicer
formatted slightly differently. Either way you are dicing with death. Of 
course when a human looks at the output of your program then they can make 
allowances for slight variations in output. When the output is read by another
computer program however, then we might not have this same flexibility. And
of course how are we to know who is going to process the output, it could
be a human today and another computer program tomorrow. So rather than risk
things going wrong we just stick to the specification. This way we have met 
our side of the agreement, and if something screws up then this is (hopefully)
someone else's problem (legal case, prison sentence or whatever).\\

\noindent
If you run tests on a number of mazes for a sufficiently long period of time, 
then you might notice some pattern in the directions the robot chooses.
You should find that the directions chosen by the {\bf DumboController} are 
indeed random, but that the directions are not chosen with the same 
probability\footnote{If you 
did not spot this previously then this is the difference between the {\bf DumboController}
and the {\bf RandomRobotController}.}. \\

\noindent
$\diamondsuit$ Investigate the definitions of the {\tt Math.random()} and {\tt Math.round()} methods from
the Java API. 
Using this information state the probability of the robot choosing the directions {\em left},
{\em right}, {\em ahead} and {\em behind}. \\

\noindent
{\bf Hint:} The Java APIs are an excellent source of code that tens of thousands of 
programmers draw upon every day. They provide a repository of program code 
that you can use to construct more complex code. A word of warning though; just because 
someone else has written the code, doesn't mean that when you employ it you are excused 
the task of understanding exactly what that code does. Using the {\tt Math.random()} and
{\tt Math.round()} methods does indeed allow the robot to generate random directions, but 
probably not in the same way that you first thought. 

\subsection{Exercise 2}

$\diamondsuit$
After further discussion with your customer they submit a revised specification.

\begin{quote}
`Could we have a modified robot controller that chooses directions randomly with equal
probability? The robot should still avoid crashing into walls and should still print a
log of its movements.'
\end{quote}

\noindent
You should base your solution to this exercise on the previous exercise. So make sure
that you continue working on the {\tt DumboController.java} file. You should also remember
to keep regular back-ups at this point.\\

\noindent
{\bf Hint:}
Review your answer to Exercise 1. There are many ways to accomplish this task; think
carefully about the use of the {\tt Math.random()} and {\tt Math.round()} functions.

\subsubsection{Improving performance?}

$\diamondsuit$ 
Use the same method of {\em design}, {\em build} and {\em test} and further 
modify {\tt DumboController} so that it
meets the following customer requirements: 

\begin{quote}
`The robot should only change direction if to carry on ahead would cause 
a collision. As before, when {\tt DumboController} chooses a
direction, it should select randomly from all directions which won't
cause a collision. Directions should be chosen with equal probability and a 
log should be kept of the robot's movements.'
\end{quote}

\noindent
Your solution should be building upon the work you have done previously, so 
again modify the code in the file {\tt DumboController.java}.
The modification you are required to make is again very small. These are the 
design questions which you need to consider: \\

\noindent
{\bf Design question 1:}
How will you instruct the robot  
controller to check {\tt if} there is a wall ahead before it decides to 
change direction?\\

\noindent
{\bf Design question 2:}
What should the controller do if there is a wall ahead of the robot? \\

\noindent
{\bf Design question 3:}
What should the controller do if there is not a wall ahead? \\

\noindent
{\bf Design question 4:}
Does your controller program work the very first time it is run? What is the
initial value of the variable {\tt direction}?

%\noindent
%Copy your answer from {\tt DumboController.java} to the file 
%{\tt Ex8.java}. You will 
%notice that you are building and testing your code using the file 
%{\tt DumboController.java}.
%This is intended as this allows you to develop code in one file and then 
%use the files {\tt Ex*.java} as backups,  
%for marking purposes, and as the code which you finally ship to the customer.
%Always remember to make comments as you go, as it will be an arduous task
%adding comments to all of your {\tt Ex*.java} files! Also make sure you rename
%your classes inside the files such that they compile from the {\tt Ex*.java} file correctly.

\subsubsection{Reaching the end}

You may notice your robot has some issues after implementing the last algorithm.
To correct this issue, it is decided some elements of randomness will be added
back into the robot controller. \\

\noindent $\diamondsuit$ 
Again building on your previous work, modify the controller
in {\tt DumboController.java}. This time, as
well as displaying all the previous characteristics, the robot controller 
should also choose a new direction randomly, on \emph{average} every 1 in 8 moves,
irrespective of whether there is a wall ahead or not. \\

\noindent
There are two design questions which need to be considered: \\

\noindent
{\bf Design question 1:}
How will you get your robot controller to choose a new direction on
average every 1 in 8 moves? Hint: you have given this some thought earlier. \\

\noindent
{\bf Design question 2:} 
How will you incorporate this into the existing code? Can you combine 
this new code with your existing {\tt if} statement by employing 
some boolean algebra? Hint: look at the lecture notes and your textbooks 
for information on logical operators in Java. \\

\noindent
Once you have established your design, build the new robot controller. 
Again, your solution should modify no more than two lines of the
controller program, so do not get too carried away. \\

\noindent
One example of the resulting behaviour 
will be that a robot which is travelling forwards in a straight line will
occasionally change direction. It may choose to go backwards,
continue on its original path forwards, or choose a left or a right
turn if these are available. Can this version of {\bf DumboController} be
expected to reach the target if given enough time? Test your solution on 
some suitable examples. \\

\noindent
Make a copy of your final solution to this exercise by typing \\

{\tt cp DumboController.java Ex2.java}

\subsection{Final words}

You may have decided while doing the previous exercises that there are many 
alternative solutions to these problems; you would of course be right. Which
solution you ultimately choose is a matter of taste, though later in your 
programming career you may well find that {\em in-house styles} or 
{\em programming templates} dictate the approach which you use. \\

\noindent
Think about an alternative approach to solving one of the previous exercises. 
Why is your alternative solution better and why?\\

\noindent
{\bf Important note}\\

\noindent
Before moving onto Part 2, ensure you have correctly renamed all of your files and classes ready for submission. The marking process for this module is partially automated and any compile errors will result in a loss of marks. Additionally, ensure that you haven't accidentally pasted rubbish into your files by accidentally hitting the middle click button.


